\documentclass{article}
\usepackage{amsmath,amssymb,amsthm}
\usepackage{algorithm}
\usepackage{algpseudocode}
\usepackage{bm}
\usepackage{enumitem}
\usepackage{hyperref}
\usepackage{geometry}
\geometry{margin=1in}
\newtheorem{theorem}{Theorem}
\newtheorem{lemma}{Lemma}
\newtheorem{assumption}{Assumption}
\newtheorem{corollary}{Corollary}
\newcommand{\E}{\mathbb{E}}
\newcommand{\R}{\mathbb{R}}
\newcommand{\norm}[1]{\left\lVert #1\right\rVert}
\newcommand{\prox}{\operatorname{prox}}

\begin{document}

\section*{Random Block Coordinate Descent (RBCD)}

\section*{Mathematical Formulation}
Let $f:\R^{d}\rightarrow \R$ be a convex differentiable function.
Partition the coordinate set $\{1,\dots,d\}$ into $B$ (possibly overlapping) blocks
$\mathcal{B}_{1},\dots ,\mathcal{B}_{B}$.
For a vector $w\in\R^{d}$ we denote by $w^{(b)}$ the subvector consisting of the
coordinates in block $\mathcal{B}_{b}$ and by $\nabla_{b}f(w)$ the corresponding
partial gradient.

At iteration $t=0,1,2,\dots$ the algorithm draws a block index $i_{t}\in\{1,\dots ,B\}$
according to a fixed probability distribution $p=(p_{1},\dots ,p_{B})$
with $p_{b}>0$ and $\sum_{b}p_{b}=1$.
Only the selected block is updated:
\[
w_{t+1}^{(i_{t})}=w_{t}^{(i_{t})}-\eta_{t}\,\nabla_{i_{t}}f(w_{t}),\qquad
w_{t+1}^{(b)}=w_{t}^{(b)}\;\;\forall\,b\neq i_{t},
\tag{1}
\]
where $\eta_{t}>0$ is the stepsize (possibly constant).

\section*{Pseudocode}
\begin{algorithm}[H]
\caption{Random Block Coordinate Descent}
\begin{algorithmic}[1]
\Require Initial point $w_{0}\in\R^{d}$, stepsize schedule $\{\eta_{t}\}_{t\ge0}$,
block partition $\{\mathcal{B}_{b}\}_{b=1}^{B}$, sampling distribution $p$.
\For{$t=0,1,\dots,T-1$}
    \State Sample $i_{t}\sim p$.
    \State Compute partial gradient $\displaystyle g_{t}= \nabla_{i_{t}}f(w_{t})$.
    \State Update block $i_{t}$:
    \[
    w_{t+1}^{(i_{t})}= w_{t}^{(i_{t})}-\eta_{t} g_{t}.
    \]
    \State Keep all other blocks unchanged.
\EndFor
\State \Return $w_{T}$ (or the average $\bar w_{T}= \frac1T\sum_{t=0}^{T-1}w_{t}$).
\end{algorithmic}
\end{algorithm}

\section*{Dry Run Example}
Consider the scalar (single‑block) function
\[
f(w)=(w-3)^{2},\qquad w\in\R .
\]
Here $B=1$, $p_{1}=1$, and the gradient is $\nabla f(w)=2(w-3)$.
Choose $w_{0}=0$, constant stepsize $\eta=0.1$, and run three iterations.

\begin{center}
\begin{tabular}{c|c|c|c}
$t$ & $w_{t}$ & $\nabla f(w_{t})$ & $f(w_{t})$ \\ \hline
0 & $0$ & $2(0-3)=-6$ & $9$ \\ 
1 & $0-0.1(-6)=0.6$ & $2(0.6-3)=-4.8$ & $(0.6-3)^{2}=5.76$ \\ 
2 & $0.6-0.1(-4.8)=1.08$ & $2(1.08-3)=-3.84$ & $(1.08-3)^{2}=3.6864$ \\ 
3 & $1.08-0.1(-3.84)=1.464$ & $2(1.464-3)=-3.072$ & $(1.464-3)^{2}=2.3700$ 
\end{tabular}
\end{center}
The objective value decreases at each step, illustrating the mechanics of (1).

\newpage

\section*{Assumptions}
\begin{assumption}[Convexity]\label{ass:convex}
$f$ is convex, i.e. for all $u,v\in\R^{d}$,
\[
f(v)\ge f(u)+\langle \nabla f(u),v-u\rangle .
\]
\end{assumption}

\begin{assumption}[Block‑coordinate $L$‑smoothness]\label{ass:smooth}
For each block $b$ there exists $L_{b}>0$ such that for all $w,\tilde w\in\R^{d}$
\[
f\bigl(w^{(1)},\dots ,w^{(b-1)},\tilde w^{(b)},w^{(b+1)},\dots ,w^{(B)}\bigr)
\le
f(w)+\langle \nabla_{b}f(w),\tilde w^{(b)}-w^{(b)}\rangle
+\frac{L_{b}}{2}\norm{\tilde w^{(b)}-w^{(b)}}^{2}.
\tag{2}
\]
Define $L_{\max}\triangleq \max_{b}L_{b}$.
\end{assumption}

\begin{assumption}[Bounded variance of block gradients]\label{ass:bounded}
For every block $b$ and all $w$,
\[
\norm{\nabla_{b}f(w)}^{2}\le G_{b}^{2},
\qquad G\triangleq \max_{b}G_{b}<\infty .
\]
\end{assumption}

\begin{assumption}[Stepsize]\label{ass:stepsize}
The deterministic stepsize satisfies 
\[
0<\eta_{t}\le \frac{1}{L_{\max}}\;\;\forall t .
\]
\end{assumption}

\section*{Main Convergence Theorem}
\begin{theorem}[Expected suboptimality of RBCD]\label{thm:main}
Under Assumptions~\ref{ass:convex}–\ref{ass:stepsize} and a constant stepsize
$\eta_{t}\equiv\eta\le 1/L_{\max}$,
the iterates generated by \eqref{1} satisfy for any $T\ge1$,
\[
\E\!\left[ f(\bar w_{T})\right]-f(w^{\star})
\;\le\;
\frac{\norm{w_{0}-w^{\star}}^{2}}{2\eta\,T},
\qquad 
\bar w_{T}\triangleq\frac{1}{T}\sum_{t=0}^{T-1}w_{t},
\tag{3}
\]
where $w^{\star}\in\arg\min f$.
Consequently, the expected optimality gap decays as $O(1/T)$.
\end{theorem}

\section*{Theoretical Properties}
\begin{itemize}[leftmargin=*]
\item \textbf{Stability.} Because $\eta\le 1/L_{\max}$, each block update is a
non‑expansive mapping in the sense of (2); the iterates remain bounded.
\item \textbf{Hyper‑parameter sensitivity.} The bound \eqref{3} is monotone in
$\eta$ for $0<\eta\le 1/L_{\max}$; the optimal choice in the bound is
$\eta=1/L_{\max}$, yielding the smallest constant.
\item \textbf{Comparison to full‑gradient GD.} When $B=1$ (single block) the bound
coincides with the classical $O(1/T)$ rate of gradient descent.
When $B=d$ and each block contains one coordinate, the method reduces to
randomized coordinate descent; the rate matches the well‑known
$\frac{d}{2\eta T}$ bound after scaling the stepsize by $d$.
\end{itemize}

\section*{Discussion}
The theorem shows that random block selection does not deteriorate the
asymptotic $1/T$ rate for convex smooth objectives, provided the stepsize is
chosen according to the worst block smoothness constant $L_{\max}$.  The
expectation is taken with respect to the random block sequence
$\{i_{t}\}_{t\ge0}$ only; no stochastic gradient noise is required.

\newpage

\section*{Preliminaries (Definitions and Lemmas)}
\begin{lemma}[Three‑point inequality for block smoothness]\label{lem:3pt}
Under Assumption~\ref{ass:smooth}, for any block $b$ and any $w,\tilde w$,
\[
f(w)-f\bigl(w^{(b)}\!\!\to\!\tilde w^{(b)}\bigr)
\ge
\langle \nabla_{b}f(w),w^{(b)}-\tilde w^{(b)}\rangle
-\frac{L_{b}}{2}\norm{w^{(b)}-\tilde w^{(b)}}^{2}.
\tag{4}
\]
\end{lemma}
\begin{proof}
Inequality \eqref{2} applied to $\tilde w^{(b)}$ and then rearranged
gives \eqref{4}. \hfill\qedsymbol
\end{proof}

\section*{Main Proof (Step‑by‑Step Derivation)}
We prove Theorem~\ref{thm:main}. Throughout the proof, $w^{\star}$ denotes any
minimizer of $f$.

\noindent\textbf{Notation.}
Let $\mathcal{F}_{t}$ be the $\sigma$‑algebra generated by the random blocks
$\{i_{0},\dots ,i_{t-1}\}$.  Conditional expectation $\E[\cdot|\mathcal{F}_{t}]$
is taken over the random choice $i_{t}$ only.

\medskip
\noindent\textbf{Step 1.}  Expand the squared distance after one update.
\begin{align}
\norm{w_{t+1}-w^{\star}}^{2}
&= \norm{w_{t}-\eta\nabla_{i_{t}}f(w_{t})-w^{\star}}^{2} \nonumber\\
&= \norm{w_{t}-w^{\star}}^{2}
-2\eta\langle \nabla_{i_{t}}f(w_{t}),\, w_{t}^{(i_{t})}-w^{\star (i_{t})}\rangle
+\eta^{2}\norm{\nabla_{i_{t}}f(w_{t})}^{2}. \tag{5}
\end{align}
[Step 1] uses the identity $\|a+b\|^{2}=\|a\|^{2}+2\langle a,b\rangle+\|b\|^{2}$
and the fact that only block $i_{t}$ changes.

\medskip
\noindent\textbf{Step 2.}  Take conditional expectation w.r.t.\ $\mathcal{F}_{t}$.
Because $i_{t}$ is drawn independently of $\mathcal{F}_{t}$,
\[
\E\!\left[\langle \nabla_{i_{t}}f(w_{t}),\, w_{t}^{(i_{t})}-w^{\star (i_{t})}\rangle
\Big|\mathcal{F}_{t}\right]
= \sum_{b=1}^{B}p_{b}\,
\langle \nabla_{b}f(w_{t}),\, w_{t}^{(b)}-w^{\star (b)}\rangle . \tag{6}
\]
Similarly,
\[
\E\!\left[\norm{\nabla_{i_{t}}f(w_{t})}^{2}\Big|\mathcal{F}_{t}\right]
= \sum_{b=1}^{B}p_{b}\,\norm{\nabla_{b}f(w_{t})}^{2}. \tag{7}
\]

\medskip
\noindent\textbf{Step 3.}  Apply Lemma~\ref{lem:3pt} with $\tilde w^{(b)}=w^{\star (b)}$.
For each block $b$,
\[
\langle \nabla_{b}f(w_{t}),\, w_{t}^{(b)}-w^{\star (b)}\rangle
\ge
f(w_{t})-f(w^{\star})-\frac{L_{b}}{2}\norm{w_{t}^{(b)}-w^{\star (b)}}^{2}.
\tag{8}
\]

\medskip
\noindent\textbf{Step 4.}  Plug \eqref{8} into the conditional expectation of
\eqref{5}. Using \eqref{6} and \eqref{7} we obtain
\begin{align}
\E\!\left[\norm{w_{t+1}-w^{\star}}^{2}\Big|\mathcal{F}_{t}\right]
&\le \norm{w_{t}-w^{\star}}^{2}
-2\eta\!\sum_{b}p_{b}\!\Bigl[f(w_{t})-f(w^{\star})
-\frac{L_{b}}{2}\norm{w_{t}^{(b)}-w^{\star (b)}}^{2}\Bigr]\nonumber\\
&\qquad +\eta^{2}\!\sum_{b}p_{b}\norm{\nabla_{b}f(w_{t})}^{2}. \tag{9}
\end{align}
[Step 4] uses linearity of expectation and the inequality direction from (8).

\medskip
\noindent\textbf{Step 5.}  Upper‑bound the gradient term by smoothness.
From Assumption~\ref{ass:smooth} with $\tilde w^{(b)}=w_{t}^{(b)}-\frac{1}{L_{b}}\nabla_{b}f(w_{t})$,
\[
f\bigl(w_{t}^{(b)}-\frac{1}{L_{b}}\nabla_{b}f(w_{t})\bigr)
\le
f(w_{t})-\frac{1}{2L_{b}}\norm{\nabla_{b}f(w_{t})}^{2}.
\]
Since $f(w^{\star})\le f\bigl(w_{t}^{(b)}-\frac{1}{L_{b}}\nabla_{b}f(w_{t})\bigr)$,
\[
f(w_{t})-f(w^{\star})\ge \frac{1}{2L_{b}}\norm{\nabla_{b}f(w_{t})}^{2}.
\tag{10}
\]
Rearranging gives
\[
\norm{\nabla_{b}f(w_{t})}^{2}\le 2L_{b}\bigl[f(w_{t})-f(w^{\star})\bigr].
\tag{11}
\]

\medskip
\noindent\textbf{Step 6.}  Substitute \eqref{11} into \eqref{9} and use
$L_{b}\le L_{\max}$:
\begin{align}
\E\!\left[\norm{w_{t+1}-w^{\star}}^{2}\Big|\mathcal{F}_{t}\right]
&\le \norm{w_{t}-w^{\star}}^{2}
-2\eta\!\sum_{b}p_{b}\!\Bigl[f(w_{t})-f(w^{\star})\Bigr]
+\eta^{2}\!\sum_{b}p_{b}\,2L_{b}\!\bigl[f(w_{t})-f(w^{\star})\bigr]\nonumber\\
&\qquad
+ \eta\!\sum_{b}p_{b}L_{b}\norm{w_{t}^{(b)}-w^{\star (b)}}^{2}. \tag{12}
\end{align}
Because $\sum_{b}p_{b}=1$ and $L_{b}\le L_{\max}$,
\[
\eta^{2}\sum_{b}p_{b}2L_{b}\le 2\eta^{2}L_{\max}\le 2\eta,
\quad\text{(using }\eta\le 1/L_{\max}\text{)}.
\tag{13}
\]

\medskip
\noindent\textbf{Step 7.}  Combine the coefficients of $f(w_{t})-f(w^{\star})$.
From \eqref{12} and \eqref{13},
\[
\E\!\left[\norm{w_{t+1}-w^{\star}}^{2}\Big|\mathcal{F}_{t}\right]
\le \norm{w_{t}-w^{\star}}^{2}
-2\eta\bigl[1-\eta L_{\max}\bigr]\bigl[f(w_{t})-f(w^{\star})\bigr]
+\eta L_{\max}\norm{w_{t}-w^{\star}}^{2}.
\tag{14}
\]
Since $\eta\le 1/L_{\max}$, the factor $1-\eta L_{\max}\ge0$.
Rearranging (14) yields
\[
2\eta\bigl[1-\eta L_{\max}\bigr]\bigl[f(w_{t})-f(w^{\star})\bigr]
\le \norm{w_{t}-w^{\star}}^{2}
-\E\!\left[\norm{w_{t+1}-w^{\star}}^{2}\Big|\mathcal{F}_{t}\right].
\tag{15}
\]

\medskip
\noindent\textbf{Step 8.}  Take total expectation and sum over $t=0,\dots,T-1$.
Denote $\Delta_{t}\triangleq \E\!\bigl[\norm{w_{t}-w^{\star}}^{2}\bigr]$.
Summing \eqref{15} gives
\[
2\eta\bigl[1-\eta L_{\max}\bigr]\sum_{t=0}^{T-1}
\E\!\bigl[f(w_{t})-f(w^{\star})\bigr]
\le \Delta_{0}-\Delta_{T}\le \Delta_{0}.
\tag{16}
\]

\medskip
\noindent\textbf{Step 9.}  Use convexity to relate the average iterate.
Convexity of $f$ implies
\[
f(\bar w_{T})\le \frac{1}{T}\sum_{t=0}^{T-1}f(w_{t}),
\qquad
\bar w_{T}\triangleq\frac{1}{T}\sum_{t=0}^{T-1}w_{t}.
\tag{17}
\]
Applying expectation and (16) with $\eta\le 1/L_{\max}$ (so that $1-\eta L_{\max}\ge 0$)
gives
\[
\E\!\bigl[f(\bar w_{T})-f(w^{\star})\bigr]
\le \frac{\Delta_{0}}{2\eta\bigl[1-\eta L_{\max}\bigr]\,T}
\le \frac{\norm{w_{0}-w^{\star}}^{2}}{2\eta T},
\tag{18}
\]
where the last inequality uses $1-\eta L_{\max}\le 1$.
\eqref{18} is exactly the bound \eqref{3}, completing the proof.
\hfill\qedsymbol

\section*{Rate Derivation (With Constants)}
The bound \eqref{3} can be expressed with explicit constants:
\[
\E\!\bigl[f(\bar w_{T})\bigr]-f(w^{\star})
\le
\frac{L_{\max}}{2T}\,\norm{w_{0}-w^{\star}}^{2},
\quad\text{by choosing }\eta=\frac{1}{L_{\max}}.
\]
Thus the convergence constant is proportional to the worst block smoothness
parameter.

\section*{Special Cases}
\begin{enumerate}
\item \textbf{Uniform block sampling.}
If $p_{b}=1/B$ for all $b$, the expectation in (6)–(7) becomes an average over
blocks, but the final rate \eqref{3} is unchanged.
\item \textbf{Single block ($B=1$).}
RBCD reduces to standard gradient descent; $L_{\max}=L_{1}$ and the bound
coincides with the classical $O(1/T)$ result.
\item \textbf{Coordinate descent ($B=d$, each block a scalar).}
Choosing $\eta=1/L_{\max}$ yields
\[
\E\!\bigl[f(\bar w_{T})-f(w^{\star})\bigr]\le
\frac{d}{2\,T}\max_{i}L_{i}\,\norm{w_{0}-w^{\star}}^{2},
\]
which matches the known rate for randomized coordinate descent after
absorbing the factor $d$ into the stepsize.
\end{enumerate}

\end{document}